\textbf{Proof.}
Use normalized representatives
\[
K_\kappa=
\left\{(b_1,b_0):
\lVert b_1\rVert_2^2=\lVert b_0\rVert_2^2=n,
\quad S_B-n\kappa I_2\succeq0
\right\}.
\tag{1}
\]
Assumption \(\operatorname{ass{:}kappa\mbox{-}feasibility}\) makes \(K_\kappa\) nonempty.  It is closed in the product of two radius-\(\sqrt n\) spheres and hence compact.  Proposition \(\operatorname{prop{:}kappa\mbox{-}exhaustion}\) identifies its projective quotient with \(\mathcal B_\kappa\).  Moreover, (1) gives
\[
S_B\succeq n\kappa I_2,
\qquad
\lVert S_B^{-1}\rVert_2\le (n\kappa)^{-1}
\quad(B\in K_\kappa).
\tag{2}
\]
Thus every inverse used by \(w\) exists, including on the \(\kappa\)-equality boundary.  Assumption \(\operatorname{ass{:}exposure\mbox{-}positivity}\) makes \(\Pi^{-1}\) well defined.

For \(e\in\{L,U\}\), put
\[
R_e=\Omega\Pi^{-1}Q(c_e)\Pi^{-1}\Omega,
\qquad
S(B)=B^{\top}\Omega B,
\qquad
T_e(B)=B^{\top}R_eB.
\tag{3}
\]
The exact algebraic input permits computation of rational bounds \(M_\Omega\ge\lVert\Omega\rVert_2\) and \(M_e\ge\lVert R_e\rVert_2\).  If \(B,C\in K_\kappa\) and \(d=\lVert B-C\rVert_F\), then \(\lVert B\rVert_2,\lVert C\rVert_2\le\sqrt{2n}\), whence
\[
\lVert S(B)-S(C)\rVert_2
\le 2\sqrt{2n}\,M_\Omega d,
\qquad
\lVert T_e(B)-T_e(C)\rVert_2
\le 2\sqrt{2n}\,M_e d.
\tag{4}
\]
The resolvent identity gives
\[
S(B)^{-1}-S(C)^{-1}
=S(B)^{-1}\{S(C)-S(B)\}S(C)^{-1}.
\tag{5}
\]
Combining (2)--(5), \(\lVert T_e(B)\rVert_2\le2nM_e\), and \(|\operatorname{tr}(UV)|\le2\lVert U\rVert_2\lVert V\rVert_2\) for two-by-two matrices yields
\[
\left|
n^{-2}\operatorname{tr}\{S(B)^{-1}T_e(B)\}
-n^{-2}\operatorname{tr}\{S(C)^{-1}T_e(C)\}
\right|
\le L_e d,
\tag{6}
\]
where, for example, the following effective rational upper bound is valid:
\[
L_e=
n^{-2}\left{
\frac{8\sqrt{2n}\,M_\Omega M_e}{n\kappa^2}
+\frac{4\sqrt{2n}\,M_e}{n\kappa}
\right}.
\tag{7}
\]
Increasing algebraic square-root bounds outward if necessary makes (7) rational.  The terms \(h(c_e)\) do not depend on \(B\), and the maximum of finitely many functions with common Lipschitz bound has that bound.  Hence an effectively computable positive rational \(\overline L\) satisfies
\[
|w(B)-w(C)|\le\overline L\lVert B-C\rVert_F
\quad(B,C\in K_\kappa).
\tag{8}
\]

Every normalized arm vector has a coordinate of absolute value at least \(1\), so the finitely many signed pivot charts with threshold \(n^{-1/2}\) cover the quotient.  Fix a rational \(\rho>0\) such that \(4\rho<\varepsilon\).  Uniformly bisect this finite chart cover until every resulting box \(C\) has diameter \(d_C\) satisfying
\[
d_C\le\varepsilon,
\qquad
\overline Ld_C\le\rho.
\tag{9}
\]
Only finitely many bisections are required.  By \(\operatorname{lem{:}algebraic\mbox{-}box\mbox{-}feasibility\mbox{-}cad}\), every box is then either deleted with a finite exact infeasibility certificate or supplied with a feasible point \(B_C\) having finite real-algebraic coordinates.  Exact algebraic evaluation and root isolation return rational numbers \(\underline q_C,\overline q_C\) satisfying
\[
w(B_C)-\rho\le\underline q_C\le w(B_C)
\le\overline q_C\le w(B_C)+\rho.
\tag{10}
\]
For every feasible box define
\[
\lambda_C=\underline q_C-\overline Ld_C.
\tag{11}
\]
For \(B\in C\cap K_\kappa\), equations (8), (10), and (11) imply
\[
\lambda_C
\le w(B_C)-\overline Ld_C
\le w(B).
\tag{12}
\]
Thus \(\lambda_C\) is a certified rational lower bound on the objective throughout the feasible part of \(C\).

Set
\[
L_\varepsilon=\min_C\lambda_C,
\qquad
U_\varepsilon=\min_C\overline q_C,
\tag{13}
\]
and let \(B_\varepsilon=B_C\) for a box attaining the second minimum.  The minima in (13) are over a nonempty finite family.  Compactness and continuity give a minimizer \(B^\star\in K_\kappa\) with \(w(B^\star)=v_{\mathrm{mm}}\).  Equations (9)--(12) give, for every feasible box, \(\lambda_C\ge v_{\mathrm{mm}}-2\rho\); applying (8) and (10) in a box containing \(B^\star\) gives \(\overline q_C\le v_{\mathrm{mm}}+2\rho\).  Consequently,
\[
v_{\mathrm{mm}}-2\rho
\le L_\varepsilon
\le v_{\mathrm{mm}}
\le w(B_\varepsilon)
\le U_\varepsilon
\le v_{\mathrm{mm}}+2\rho.
\tag{14}
\]
In particular, \(U_\varepsilon-L_\varepsilon\le4\rho<\varepsilon\).

Prune a feasible box only when its certified lower bound \(\lambda_C\) is strictly larger than \(U_\varepsilon\).  By (12)--(14), no box containing a global minimizer can be pruned.  Therefore the union \(\mathcal C_\varepsilon\) of the unpruned quotient-chart boxes satisfies
\[
\mathcal A_{\mathrm{mm}}\subseteq\mathcal C_\varepsilon.
\tag{15}
\]
Equation (9) gives \(h_\varepsilon\le\varepsilon\).  The incumbent is one of the CAD samples, so it is exactly feasible and has finite real-algebraic coordinates even if all feasible points satisfy the regularity constraint with equality.

The finite record consists precisely of the chart cover, CAD sign and feasibility certificates, the algebraic incumbent, outward rational objective bounds, rational norm and diameter records, pruning comparisons, and the mesh.  A local checker verifies those records; (12)--(15) then give the asserted statements about the unknown value and minimizer set.  For any sequence \(\varepsilon\downarrow0\), the certified gap and mesh are bounded above by \(\varepsilon\), and therefore both converge to zero.  The argument supplies neither a uniform subdivision-complexity bound nor a semidefinite-program reduction. \(\square\)